\section{实数}
\label{sec:ru-real-numbers}
\label{sec:real-valued-functions-bounds}

本节沿用读者熟悉的实数及其四则运算和大小关系，不讨论实数如何从有理数构造。
我们只整理后文反复使用的三样东西：描述位置与误差的区间和绝对值，保证“最小上界”存在的完备性，以及由完备性得到的十进制表示。

在实轴上，区间按端点是否包含分为几种：$[a,b]=\{x:a\leq x\leq b\}$称为闭区间，$(a,b)=\{x:a<x<b\}$称为开区间，
$[a,b)$与$(a,b]$称为半开区间；无界区间写作$[a,+\infty)$、$(-\infty,b)$等，$(-\infty,+\infty)=\mathbb R$。
这里的$\pm\infty$只是记号，表示区间在那一侧没有端点，它们不是实数。
数$x$的\term{绝对值}（absolute value）$|x|$在$x\geq0$时等于$x$，否则等于$-x$；$|x-y|$就是实轴上两点之间的距离。
于是$|x-a|<r$恰好表示$x\in(a-r,a+r)$，即$x$落在以$a$为中心、半径为$r$的开区间内。

由$-|x|\leq x\leq|x|$与$-|y|\leq y\leq|y|$相加，得$-(|x|+|y|)\leq x+y\leq|x|+|y|$，即\term{三角不等式}（triangle inequality）
\[
 |x+y|\leq|x|+|y|.
\]
取$x-y$与$y-z$代入，得到$|x-z|\leq|x-y|+|y-z|$：经过第三点绕行，距离不会变短。
把$x=(x-y)+y$代入三角不等式又得$|x|-|y|\leq|x-y|$，交换$x,y$后合并，就是反三角不等式
\[
 \bigl||x|-|y|\bigr|\leq|x-y|.
\]
它说明两数接近时，它们的绝对值也接近。

用$\hat x$近似真值$x$时，$|\hat x-x|$称为\term{绝对误差}（absolute error）。
同样的绝对误差对不同大小的真值意义不同：在$1000$上差$1$与在$2$上差$1$，后者严重得多。
因此常用相对误差$|\hat x-x|/|x|$，它把误差与真值的大小比较，所以要求参照值$x\neq0$。
例如用$3.14$近似$\pi$，绝对误差小于$0.0016$，相对误差约为$0.05\%$。

\paragraph{确界与完备性}
\label{sec:ru-supremum-completeness}
集合$A=\{1-1/n:n=1,2,\ldots\}$的元素越来越接近$1$，却没有最大的元素：任一元素$1-1/n$都小于下一个元素$1-1/(n+1)$。
$1$、$2$、$100$都不小于$A$中的每个数，而在这些数中，$1$是最小的一个。下面把这种“最小上界”说清楚。

设$A\subseteq\mathbb R$非空。若$a\leq M$对每个$a\in A$成立，就称$M$为$A$的上界；若$m\leq a$对每个$a\in A$成立，就称$m$为$A$的下界。
存在上界称为上有界，存在下界称为下有界，两者都有称为有界。界不必属于集合，也不唯一。
这正是第~\ref{sec:prelim-order-relations}节偏序集合中上界、下界的定义在$(\mathbb R,\leq)$中的情形。

\begin{definition}[上确界与下确界]
\label{def:ru-supremum}
设$A\subseteq\mathbb R$非空。若$\beta$是$A$的上界，并且任何$A$的上界$M$都满足$\beta\leq M$，就称$\beta$为$A$的\term{上确界}（supremum），记作$\sup A$。
类似地，下界中最大的一个称为\term{下确界}（infimum），记作$\inf A$。
若$\sup A\in A$，就称它为$A$的最大值，记作$\max A$；若$\inf A\in A$，就称它为最小值，记作$\min A$。
\end{definition}

上确界的第二个条件常改写为：对每个$\varepsilon>0$，都存在$a\in A$使$a>\beta-\varepsilon$。
它说的是任何比$\beta$小的数都不再是上界。下确界$\alpha$相应地满足：对每个$\varepsilon>0$，都存在$a\in A$使$a<\alpha+\varepsilon$。
上例中$\sup A=1\notin A$，所以$A$没有最大值；$\inf A=\min A=0$，因为$0=1-1/1\in A$。
可见界、确界与最值是三个层次：界有许多，确界至多一个，最值则还要求确界属于集合。

上界有无限多个，其中是否真有最小的一个，并不能从定义看出。
本书把下面的性质作为实数的基本性质使用，称为实数的\term{完备性}（completeness），也称\term{确界原理}（least-upper-bound property）：
非空且上有界的实数集合有实数上确界；非空且下有界的实数集合有实数下确界。
后一句可由前一句得到：$A$下有界时，$-A=\{-a:a\in A\}$上有界，$-\sup(-A)$就是$\inf A$。

有理数没有这一性质。集合$\{q\in\mathbb Q:q^2<2\}$在有理数中有上界$2$，它在实数中的上确界却是$\sqrt2$，而$\sqrt2$不是有理数。
于是每个有理上界都严格大于$\sqrt2$，在它与$\sqrt2$之间总还有更小的有理上界，有理数中找不到最小的上界。
实数补上了这样的“空隙”。下面三条性质都由确界原理推出，下一章证明极限存在时会直接使用它们。

第一条是\term{Archimedes性}（Archimedean property）：对任意实数$x$，都存在正整数$n>x$。
若不然，正整数集合以$x$为上界，由确界原理有上确界$s$；$s-1$不是上界，故有正整数$k>s-1$，于是正整数$k+1>s$，矛盾。
等价地，对任意$\varepsilon>0$，都有正整数$n$使$1/n<\varepsilon$：取$n>1/\varepsilon$即可。
由此还可以说明，对每个实数$y$都有唯一的整数$m$满足$m\leq y<m+1$，记作$m=\lfloor y\rfloor$，称为$y$的整数部分；
Archimedes性保证$\{k\in\mathbb Z:k\leq y\}$非空且有上界，它的最大元就是$m$。

第二条是有理数的稠密性：任意两个实数$a<b$之间都有有理数。
由Archimedes性取正整数$n$使$1/n<b-a$，再令$m=\lfloor na\rfloor+1$，则$m-1\leq na<m$，即
\[
 a<\frac mn\leq a+\frac1n<b.
\]
所以$m/n\in(a,b)$。区间无论多短，其中都有有理数；这与上一部分证明的“有理数只有可数多个”并不矛盾。

第三条是闭区间套性质：若闭区间$[a_1,b_1]\supseteq[a_2,b_2]\supseteq\cdots$逐个嵌套，则它们有公共点；
若还对每个$\varepsilon>0$都有某个$k$使$b_k-a_k<\varepsilon$，则公共点恰有一个。
证明如下：嵌套使左端点$a_k$单调不减，并且每个$b_j$都是全部左端点的上界（$k\leq j$时$a_k\leq a_j\leq b_j$，$k>j$时$a_k\leq b_k\leq b_j$）。
由确界原理，$c=\sup_ka_k$存在，并满足$a_k\leq c\leq b_k$对每个$k$成立，所以$c$是公共点。
若$x<y$都是公共点，则每个$k$都有$b_k-a_k\geq y-x$，与区间长度可以任意小矛盾。
端点必须属于区间：开区间套$(0,1)\supseteq(0,1/2)\supseteq(0,1/3)\supseteq\cdots$没有公共点，因为任何$x>0$都会在$1/k<x$时被排除。

\paragraph{十进制表示}
\label{sec:ru-decimal-representation}
上一部分证明实数不可数时（第~\ref{sec:comb-real-uncountable}节），借用了十进制表示的两个事实：每个实数都有十进制表示，
两种表示表示同一个数只发生在以全$0$与全$9$结尾的情形。闭区间套性质为这两个事实提供了依据。

先看实数怎样产生数字。设$x\in[0,1)$，令$x_k=\lfloor10^kx\rfloor/10^k$，它是把$x$截断到小数点后第$k$位所得的有限小数。
由$\lfloor10^kx\rfloor\leq10^kx<\lfloor10^kx\rfloor+1$，得到
\[
 x\in I_k=\bigl[x_k,\ x_k+10^{-k}\bigr),\qquad k=1,2,\ldots.
\]
从$I_k$进到$I_{k+1}$，是把长为$10^{-k}$的区间十等分，再选出$x$所在的那一份；所选份的序号$d_{k+1}\in\{0,1,\ldots,9\}$就是第$k+1$位数字。
这样得到的数字序列$0.d_1d_2d_3\cdots$不会从某位起全是$9$。若第$k$位以后全是$9$，则$j>k$时$x\geq x_j=x_k+10^{-k}-10^{-j}$；
由Archimedes性，$10^{-j}$可以任意小，于是$x\geq x_k+10^{-k}$，与$x\in I_k$矛盾。

反过来，任给数字序列$d_1,d_2,\ldots$，令$s_k=\sum_{j=1}^kd_j10^{-j}$，闭区间$J_k=[s_k,\ s_k+10^{-k}]$逐个嵌套，长度$10^{-k}$可以任意小。
由闭区间套性质，它们恰有一个公共点，这就是数字序列所表示的实数。
对前面由$x$产生的数字，$x$属于每个$J_k$，所以$x$正是它的数字序列所表示的数；这说明每个实数都有十进制表示。

一个数可能有两种表示。$0.5000\cdots$对应$J_k=[0.5,\ 0.5+10^{-k}]$，$0.4999\cdots$对应$J_k=[0.5-10^{-k},\ 0.5]$，两组区间的唯一公共点都是$0.5$。
一般地，设两个数字序列在第$k$位首次不同，且$d_k<d'_k$。第一个序列的区间不超出$p=s_{k-1}+(d_k+1)10^{-k}$，第二个序列的区间不低于$s_{k-1}+d'_k10^{-k}\geq p$。
两数相等时，公共值只能是$p$，这又要求$d'_k=d_k+1$，第一个序列此后全为$9$、第二个此后全为$0$。
所以表示不唯一恰好发生在全$9$与全$0$结尾的情形，这正是实数不可数的证明所需的两点。
数系怎样从自然数逐步扩充、实数怎样由有理数严格构造，本书不再展开。
